\chapter{Introducción}
\label{chap:intro}

Cuando se habla del impacto de la ingeniería informática en la sociedad, el
discurso predominante suele orbitar en torno al crecimiento económico: la
digitalización de los mercados, la automatización industrial o el surgimiento
de nuevos modelos de negocio basados en datos. Esta perspectiva no es errónea,
pero es incompleta. La ingeniería del software y, en particular, la
\textbf{inteligencia artificial} (\ac{IA}), son también herramientas de
progreso humano: a pesar de las críticas legítimas sobre su impacto social,
estas tecnologías pueden ayudar a reducir desigualdades, ampliar el acceso a
servicios especializados y poner en manos de colectivos vulnerables recursos
que de otro modo les serían inaccesibles.

La comunicación es una de las características más fundamentales y propias del ser humano. 
A través de ella construimos relaciones, expresamos necesidades, compartimos conocimiento y
participamos en la vida social. Sin embargo, para un porcentaje significativo de
la población (personas con parálisis cerebral, trastornos del espectro
autista, enfermedades neuromusculares degenerativas o lesiones adquiridas del
sistema nervioso) la comunicación oral resulta difícil, extremadamente
limitada o, incluso, imposible. Estas personas recurren a los \textbf{\ac{SAAC}}: tableros de pictogramas,
dispositivos generadores de voz o aplicaciones que les permiten construir
mensajes combinando símbolos visuales. Los \ac{SAAC} no son un sustituto pobre
del lenguaje oral; son, en muchos casos, la única vía real de acceso a la
comunicación, y con ella, a la autonomía, a la educación y a la participación
en comunidad.

La logopedia y la atención temprana disponen de herramientas para evaluar el
desarrollo lingüístico de estos usuarios, identificar sus necesidades
específicas y diseñar intervenciones personalizadas. No obstante, entre las
herramientas disponibles y su uso sistemático en la práctica clínica cotidiana
existe una brecha que tiene un nombre concreto: \textbf{tiempo}. El análisis riguroso
de una muestra de lenguaje asistido puede consumir mucho tiempo de
trabajo experto por cada sesión. Ante esa carga, muchos profesionales se ven 
forzados a priorizar la intervención directa sobre la evaluación detallada, 
o a espaciar las valoraciones hasta hacerlas poco representativas de la 
evolución real del usuario.

Es en este punto donde la ingeniería informática tiene algo que aportar que va
más allá de la eficiencia. Automatizar una parte de ese análisis no significa
reemplazar al logopeda ni reducir la complejidad clínica a un algoritmo, sino
\textit{liberar} al profesional de las tareas más mecánicas y repetitivas
(la \textbf{transcripción}, el \textbf{conteo de morfemas}, el
\textbf{cómputo de métricas estandarizadas}) para que pueda dedicar su tiempo
y su criterio a lo que ningún sistema informático puede hacer: interpretar,
contextualizar, relacionarse con el usuario y con su familia, y tomar
decisiones terapéuticas fundadas en una evaluación más frecuente y completa.
Este trabajo busca precisamente ese objetivo, poniendo las técnicas del
Procesamiento del Lenguaje Natural (\ac{PLN}), el reconocimiento automático
del habla y el aprendizaje automático al servicio de un colectivo que las
necesita: las personas que se comunican mediante \ac{SAAC} y los profesionales
que las acompañan.

\section{Contexto y motivación}

El \textbf{análisis de muestras de lenguaje} (\emph{Language Sample Analysis}, LSA)
constituye una metodología ampliamente utilizada en logopedia para caracterizar las
habilidades comunicativas expresivas de una persona. Consiste en recoger una muestra representativa de la 
producción lingüística del usuario en contextos naturales, transcribirla y analizarla 
sistemáticamente en múltiples dimensiones: el vocabulario empleado, la morfología de 
las formas producidas, la complejidad de las estructuras gramaticales y la organización
sintáctica de los enunciados. El \ac{LSA} refleja el lenguaje real del usuario, no su 
capacidad de responder a un estímulo artificial, y permite un seguimiento longitudinal
genuino del progreso terapéutico \cite{lsa_computer_programs}.

Para los usuarios de \ac{SAAC}, el LSA presenta desafíos específicos que las
herramientas existentes, fundamentalmente \ac{SALT} (\emph{Systematic Analysis of
Language Transcripts}) y CLAN/CHILDES, no estaban diseñadas para abordar
\cite{soto2022paalss_tutorial}. Aunque \ac{SALT} dispone de recursos en español,
estos no han sido adaptados al lenguaje asistido por \ac{SAAC}. El lenguaje asistido no es simplemente habla oral 
reducida: los pictogramas representan mayoritariamente formas de cita sin flexión 
morfológica, los enunciados son más cortos y telegráficos, el mensaje puede construirse 
combinando símbolos con gestos y mirada, y la producción está condicionada por el 
vocabulario disponible en el dispositivo \cite{nlp_aac_application}. Aplicar métricas 
de habla oral a este tipo de producción introduce sesgos que infravaloran sistemáticamente la
competencia comunicativa real del usuario \cite{soto2022paalss_tutorial}.

La lingüista \textbf{Gloria Soto} publicó el \ac{PAALSS} (\emph{Protocol for the
Analysis of Aided Language Samples in Spanish}), el primer protocolo diseñado 
específicamente para el análisis de muestras de lenguaje asistido en español. 
El \ac{PAALSS} evalúa cuatro dimensiones: el vocabulario producido, la morfología 
gramatical, la complejidad de los enunciados medida mediante la Longitud Media del 
Enunciado (\ac{LME}), y la sintaxis expresada a través del orden canónico sujeto-verbo-objeto
\cite{soto2022paalss_tutorial}. Su aplicación requiere al menos 50 enunciados recogidos 
en contextos naturales de interacción, que un especialista en \ac{SAAC} debe transcribir, 
segmentar y codificar manualmente \cite{soto2024paalss_pilot}. El coste temporal de este proceso, que puede
alcanzar varias horas por cada sesión, es la principal barrera para su adopción sistemática 
en la práctica clínica y educativa.

Por otro lado, los avances recientes en \textit{\ac{PLN}} y en modelos de lenguaje
preentrenados ofrecen, por primera vez, herramientas con la madurez técnica
suficiente para abordar este problema de forma realista: el reconocimiento
automático del habla con modelos locales como \emph{Whisper}
\cite{lsa_swiss_benchmark}, la diarización de hablantes
\cite{whisper_diarization_realtime}, el análisis morfosintáctico basado en
\emph{transformers} \cite{lsa_automation_jslhr} y el aprendizaje con pocos ejemplos 
en dominios clínicos de datos escasos \cite{fewshot_clinical_nlp} conforman un 
conjunto de bloques tecnológicos que, combinados adecuadamente, pueden constituir 
un sistema capaz de asistir el análisis manual del \ac{PAALSS}.

De esta forma, el presente trabajo propone el diseño y la evaluación de un
sistema de \ac{IA} para la \textbf{automatización parcial} del protocolo \ac{PAALSS},
con el objetivo de reducir la carga de trabajo del clínico y facilitar la
aplicación sistemática del protocolo en entornos donde los recursos de tiempo
son limitados. El sistema respeta en todo momento las restricciones de
privacidad que impone el trabajo con datos de menores con discapacidad:
todo el procesamiento se realiza localmente, sin dependencia de servicios
externos en la nube \cite{lsa_swiss_benchmark}.

% TODO: Ampliar con estadísticas sobre prevalencia de usuarios SAAC
% y contexto clínico/educativo en España.

\section{Planteamiento del problema}

Dado un vídeo de una sesión de interacción entre un usuario de \ac{SAAC} y su
interlocutor, se desea obtener automáticamente los valores de las cuatro dimensiones del
protocolo \ac{PAALSS}: \textbf{vocabulario}, \textbf{morfología}, 
\textbf{complejidad gramatical} y \textbf{sintaxis}. Todo esto sin que el clínico 
tenga que realizar manualmente la transcripción, segmentación y análisis lingüístico de
los enunciados. El sistema recibe como entrada una grabación de vídeo y debe producir 
como salida un formulario con las métricas cuantitativas que el protocolo 
exige para cada una de esas dimensiones.

La naturaleza del lenguaje asistido introduce dificultades que no tienen
equivalente en los sistemas de análisis de lenguaje oral convencionales 
\cite{lsa_computer_programs}. El usuario de \ac{SAAC} construye sus mensajes combinando 
selecciones de pictogramas con vocalizaciones, gestos, mirada dirigida y, en ocasiones, 
movimientos corporales. Un sistema centrado únicamente en el canal de audio captaría en 
el mejor de los casos las vocalizaciones del usuario y del interlocutor, pero perdería 
la contribución lingüística más relevante para el análisis del protocolo: la secuencia de
pictogramas seleccionados en el tablero de comunicación. Esta asimetría entre
el canal de información primario para el análisis y el canal explotado por los
sistemas de transcripción convencionales es la primera fuente de complejidad
técnica del problema.

En segundo lugar, la morfología de las formas producidas por los usuarios de
\ac{SAAC} no es directamente observable en la señal. El análisis morfológico del
\ac{PAALSS} requiere determinar si el usuario ha producido un participio o un
plural, algo que en el habla oral se manifiesta mediante marcas fonológicas
explícitas pero que en el lenguaje asistido debe inferirse del contexto
comunicativo y de las vocalizaciones parciales que puedan acompañar la selección 
del pictograma \cite{soto2022paalss_tutorial}. Esta inferencia exige una representación
multimodal del enunciado que los modelos de \ac{PLN} convencionales, entrenados
sobre texto escrito o habla oral, no están preparados para explotar.

En tercer lugar, la escasez de datos anotados en este dominio es estructural.
El \ac{PAALSS} es un protocolo reciente y el número de sesiones anotadas según sus 
criterios en español es reducido. Los sistemas basados en \emph{transformers} que han 
obtenido resultados destacables en tareas de análisis del lenguaje oral no pueden 
reentrenarse en este dominio por falta de datos suficientes \cite{lsa_automation_jslhr}. 
El sistema debe por tanto combinar enfoques basados en reglas lingüísticas explícitas, 
análisis morfosintáctico con herramientas de \ac{PLN} en español y aprendizaje
automático con muy pocos ejemplos para las categorías más ambiguas
\cite{fewshot_clinical_nlp}.

Por otro lado, la solución técnica está acotada por la restricción de
privacidad ya introducida: el cumplimiento estricto del \textbf{Reglamento
General de Protección de Datos} (RGPD) limita el catálogo de herramientas
disponibles a modelos \emph{Open-Source} de inferencia local, como
\emph{Whisper}, con la ventaja adicional de que elimina la dependencia de 
proveedores externos y garantiza la disponibilidad del sistema
en entornos clínicos con conectividad y recursos limitados.

Cabe señalar, finalmente, lo que este trabajo no pretende resolver. El sistema
propuesto no realiza diagnósticos clínicos ni sustituye el criterio del
especialista en \ac{SAAC}. El alcance se circunscribe al español y a los formatos 
de sesión compatibles con el \ac{PAALSS}: sesiones de lectura compartida o juego 
con un interlocutor, grabadas en vídeo. La evaluación de la precisión del sistema se 
plantea mediante la comparación de sus salidas con las anotaciones expertas de dos sesiones 
reales del protocolo\cite{soto2024paalss_pilot}.

\section{Competencias adquiridas}

En cuanto a las competencias de \emph{conocimiento}, el procesamiento íntegramente
local que exigen las sesiones de una menor con discapacidad aplica la interpretación
de la legislación y los aspectos éticos del tratamiento de datos personales sensibles
(CN01), mientras que el conjunto de técnicas de \ac{IA} empleadas (reconocimiento del
habla, diarización, modelos de lenguaje preentrenados, clasificación \emph{few-shot} y
generación aumentada por recuperación) aplica la comprensión de sus fundamentos (CN02).

De las competencias de \emph{habilidad}, el \emph{pipeline} de transcripción,
diarización y fusión en Turnos Comunicativos pone en práctica la construcción de
sistemas \emph{software} basados en aprendizaje automático (HA01); la elección razonada
entre reglas lingüísticas, \ac{PLN} convencional y clasificadores \emph{few-shot} para
cada categoría morfológica ejercita la evaluación de la técnica más adecuada para cada
subproblema (HA02); y la escasez de sesiones anotadas según el protocolo \ac{PAALSS}
exige recopilar y procesar datos de calidad bajo criterios éticos y legales (HA03).

De las competencias \emph{generales}, formalizar las cuatro dimensiones del protocolo en
métricas computables aplica representación del conocimiento y aprendizaje automático
para la decisión clínica (CP01); estructurar el trabajo en cuatro objetivos específicos
ejercita la planificación de proyectos de \ac{IA} (CP02); y la identificación de
limitaciones de sesgo, trazabilidad y privacidad, junto con la renuncia explícita a toda
función diagnóstica, ejemplifica la detección de debilidades normativas y su mitigación
(CP04).


\section{Estructura del documento}

Esta memoria está compuesta por una serie de capítulos en los que se definen distintos 
aspectos a resaltar.

\begin{enumerate}
    \item \emph{Antecedentes}. Recoge la revisión bibliográfica y el estado del 
    arte en los ámbitos relevantes.
    \item \emph{Objetivos}. Presenta el objetivo general, los objetivos específicos 
    y el alcance del trabajo.
    \item \emph{Método}. Describe la metodología seguida y el diseño del sistema propuesto.
    \item \emph{Resultados}. En este capítulo es donde se exponen todo lo que se ha 
    realizado durante el desarrollo, respetando la metodología definida anteriormente. 
    Presenta los resultados obtenidos y su evaluación.
    \item \emph{Conclusiones}. Resume las conclusiones del trabajo y propone líneas de 
    trabajo futuro.
\end{enumerate}